
\label{sec:related_work}

This chapter surveys the four main research areas that collectively form the foundation of the proposed fall detection framework: thermal infrared human detection (Section~\ref{sec:rw_thermal}), human pose estimation (Section~\ref{sec:rw_pose}), skeleton-based action recognition (Section~\ref{sec:rw_action}), and vision-based fall detection systems (Section~\ref{sec:rw_fall}). For each area, this chapter traces the evolution from classical approaches to modern deep learning methods, identifies remaining challenges, and positions the contributions of this thesis within the broader literature.

\section{Thermal Infrared Human Detection}
\label{sec:rw_thermal}

\subsection*{Advantages of Thermal Imaging}

Thermal infrared imaging has gained substantial attention for human-centric perception tasks in recent years, driven by three fundamental advantages over conventional visible-spectrum cameras. First, thermal sensors capture radiated body heat rather than reflected light, which renders the resulting imagery inherently privacy-preserving: facial features, clothing textures, and other personally identifiable visual cues are largely suppressed, making it practically impossible to identify individuals from the thermal signature alone~\cite{gade2014thermal}. This property is particularly important in the context of indoor elderly monitoring, where compliance with data protection regulations such as the European General Data Protection Regulation (GDPR) imposes strict constraints on the collection and processing of biometric data. Second, because thermal emission is independent of ambient illumination, detection performance remains stable across complete darkness, strong backlighting, and varying indoor lighting conditions---a critical requirement for systems that must operate around the clock without interruption. Third, the high thermal contrast between the human body (core temperature approximately 37\textdegree{}C) and typical indoor backgrounds (20--25\textdegree{}C) provides a strong detection cue even at very low spatial resolutions. This contrast is the fundamental physical principle exploited by thermopile sensor arrays, such as the Heimann HTPA 60$\times$40~\cite{heimann2024htpa}, which produce ultra-low-resolution thermal images that maximize privacy while retaining sufficient information for human body analysis.

Gade and Moeslund~\cite{gade2014thermal} provided a comprehensive survey of thermal camera applications across domains including surveillance, medical imaging, and industrial inspection, noting that the unique characteristics of thermal imagery---particularly the absence of texture and color---require specialized processing pipelines distinct from those developed for visible-spectrum data. Early studies further demonstrated that meaningful behavioral information can be extracted from low-resolution infrared sensors even when spatial resolution is insufficient to reveal personal identity.

\subsection*{Classical and Deep Learning Detection Methods}

Early thermal human detection methods adopted classical machine learning pipelines combining handcrafted feature descriptors with discriminative classifiers. The Histogram of Oriented Gradients (HOG) descriptor introduced by Dalal and Triggs~\cite{dalal2005hog} became a foundational representation for pedestrian detection, capturing local gradient structure in a way that generalizes across appearance variation. When paired with a Support Vector Machine (SVM) classifier, HOG-based detectors achieved reasonable performance on thermal imagery but remained limited in their ability to handle diverse body poses, partial occlusions, and the scale variations inherent in practical monitoring scenarios.

The release of large-scale multispectral benchmarks substantially accelerated research in thermal detection. The KAIST Multispectral Pedestrian Detection dataset~\cite{hwang2015kaist} provided precisely aligned RGB--thermal image pairs captured across diverse urban scenes during both daytime and nighttime conditions, enabling systematic study of cross-modal fusion strategies and demonstrating that thermal imagery dramatically improves detection robustness under poor illumination. The LLVIP dataset~\cite{jia2021llvip} further extended this line of work by offering high-quality paired visible--infrared images specifically targeting low-light scenarios, where the performance gap between thermal and visible modalities is most pronounced. These benchmarks collectively established thermal human detection as a mature research direction with well-characterized advantages over single-modality visible-spectrum approaches.

With the advent of deep learning, convolutional neural network (CNN)-based object detectors rapidly surpassed classical methods. Two-stage detectors such as Faster R-CNN~\cite{ren2015faster}, which decouples region proposal generation from classification, demonstrated strong accuracy on thermal inputs but introduced computational overhead that limits real-time deployment. Single-stage detectors including YOLO~\cite{redmon2016yolo} and SSD~\cite{liu2016ssd} eliminated the proposal stage entirely, directly predicting bounding boxes and class scores from dense feature maps, achieving substantial improvements in inference speed at a modest accuracy cost. Devaguptapu~\emph{et al.}~\cite{devaguptapu2019borrow} specifically addressed the challenge of limited thermal training data by proposing a pseudo multi-modal approach that transfers representations learned from large-scale visible-spectrum datasets to thermal detectors, demonstrating that cross-domain knowledge transfer can effectively bootstrap thermal detection performance.

\subsection*{Real-Time Detection and the Proposed Approach}

Among recent architectures designed explicitly for real-time applications, RTMDet~\cite{lyu2022rtmdet} achieves an excellent speed--accuracy trade-off through several design innovations. It employs a CSPNeXt-based backbone with a Path Aggregation Feature Pyramid Network (PAFPN)~\cite{liu2018path} neck that enhances multi-scale feature fusion through bidirectional information flow. The detection head uses an anchor-free formulation that eliminates the need for pre-defined anchor boxes, simplifying the pipeline and improving generalization to objects of varying aspect ratios---a property particularly beneficial for detecting humans in the diverse poses encountered during fall events (e.g., standing, sitting, lying down). Training is supervised by Quality Focal Loss~\cite{li2020generalized}, which jointly models classification confidence and localization quality, and GIoU loss~\cite{rezatofighi2019generalized} for bounding box regression.

In the proposed framework, RTMDet is adopted as the detection architecture, but with the original CSPNeXt backbone replaced by MobileViT-S~\cite{mehta2022mobilevit}, a lightweight hybrid architecture that interleaves depthwise separable convolutions (inherited from MobileNetV2~\cite{sandler2018mobilenetv2}) with vision transformer blocks. MobileViT-S combines the local feature extraction efficiency of convolutional operations with the global contextual modeling capacity of self-attention, achieving a favorable accuracy--latency trade-off on mobile and edge hardware. The deployment target of the system---an FPGA-accelerated edge device running at $>$20~FPS---necessitates this lightweight backbone choice, as efficient neural network deployment on resource-constrained hardware remains a significant engineering challenge~\cite{sze2017efficient}.

\section{Human Pose Estimation}
\label{sec:rw_pose}

\subsection*{Top-Down versus Bottom-Up Paradigms}

Human pose estimation aims to localize anatomical keypoints of the human body from images and constitutes a fundamental building block for higher-level tasks such as action recognition, person re-identification, and human-computer interaction. Existing methods are broadly classified into two paradigms. Top-down methods~\cite{xiao2018simple, sun2019hrnet} first detect individual persons via a bounding box detector and subsequently estimate keypoints within each cropped and resized region of interest (RoI). Bottom-up approaches~\cite{cao2017realtime} detect all keypoints in the image simultaneously and group them into individual skeletons through associative embedding or part affinity fields. Cao~\emph{et al.}~\cite{cao2017realtime} demonstrated that bottom-up methods can achieve real-time multi-person pose estimation by predicting part affinity fields that encode limb associations, but their accuracy on individual instances tends to be lower than that of top-down methods, particularly when persons are at different scales or are partially occluded.

Top-down methods generally achieve higher per-instance accuracy and integrate naturally with detection-based pipelines, as the detector provides a well-defined spatial context (the bounding box) within which keypoint estimation operates. The proposed framework adopts the top-down paradigm, applying keypoint estimation to RoI-pooled features~\cite{he2017mask} extracted from each detected bounding box.

\subsection*{Heatmap-Based Methods}

Within the top-down framework, heatmap-based methods have long dominated the field. The core idea is to predict, for each keypoint, a 2D spatial probability map (heatmap) in which the peak value indicates the most likely keypoint location. The Stacked Hourglass Network~\cite{newell2016stacked} pioneered this approach by introducing repeated bottom-up and top-down processing modules with intermediate supervision, enabling the network to refine predictions progressively at multiple scales. The architecture's symmetrical encoder-decoder design with skip connections became a template for subsequent work.

SimpleBaseline~\cite{xiao2018simple} demonstrated a surprising result: a straightforward architecture comprising a ResNet backbone followed by just three deconvolutional layers could achieve accuracy competitive with far more complex designs, suggesting that the heatmap output representation itself---rather than architectural novelty---is the key driver of performance. This finding underscored the effectiveness and robustness of the Gaussian heatmap formulation as a supervision target.

HRNet (High-Resolution Network)~\cite{sun2019hrnet} further advanced the state of the art by fundamentally rethinking the backbone architecture. Rather than encoding the input to a low-resolution representation and then decoding back to high resolution (as in hourglass-style networks), HRNet maintains high-resolution feature representations in parallel throughout the network, with repeated multi-scale feature fusion. This design avoids the information loss inherent in spatial downsampling and upsampling, yielding more spatially precise heatmap predictions. HRNet achieved state-of-the-art results on the COCO keypoint benchmark~\cite{lin2014microsoft} and became one of the most widely used backbones for pose estimation.

More recently, ViTPose~\cite{xu2022vitpose} demonstrated that plain vision transformers, pre-trained on large-scale image classification datasets, can serve as highly effective backbones for pose estimation without any pose-specific architectural modifications. The global receptive field of self-attention enables ViTPose to capture long-range spatial dependencies between distant body parts, which is particularly beneficial for estimating the positions of limbs relative to the torso.

\subsection*{Limitations of 2D Heatmaps and Coordinate Refinement}

Despite the widespread success of heatmap-based methods, 2D heatmap representations suffer from an inherent quantization error: the spatial resolution of the output heatmap is typically $4\times$ to $8\times$ lower than the input image, and the standard argmax decoding step discretizes predicted coordinates to integer grid locations within this coarse spatial grid. For high-resolution inputs, this quantization may be acceptable, but for the ultra-low-resolution thermal images targeted in this thesis ($40 \times 114$ pixels), even small quantization errors can represent a significant fraction of the person's spatial extent.

Several approaches have been proposed to mitigate this limitation. DARK Pose (Distribution-Aware Coordinate Representation)~\cite{zhang2020distribution} introduced a principled sub-pixel refinement technique based on Taylor expansion of the log-likelihood around the heatmap peak, achieving measurable improvements in localization accuracy without modifying the network architecture. The integral regression (soft-argmax) approach, independently proposed by Sun~\emph{et al.}~\cite{sun2018integral} and Nibali~\emph{et al.}~\cite{nibali2018numerical}, replaced the non-differentiable argmax operation with a differentiable spatial expectation computed over the softmax-normalized heatmap. This formulation enables end-to-end gradient flow from coordinate-level supervision through the heatmap representation, combining the spatial structure of heatmaps with the precision of direct coordinate regression.

\subsection*{Direct Regression and SimDR}

At the opposite end of the design spectrum, direct coordinate regression methods predict keypoint positions directly from feature vectors without any intermediate spatial representation. DeepPose~\cite{toshev2014deeppose}, the pioneering work in this category, employed a cascade of CNN regressors to iteratively refine keypoint predictions. While conceptually simple and computationally efficient, regression-based methods have historically underperformed heatmap approaches due to the difficulty of learning the highly nonlinear mapping from high-dimensional appearance features to low-dimensional coordinate outputs. Li~\emph{et al.}~\cite{li2021human} proposed Residual Log-likelihood Estimation (RLE) to address the fundamental ill-posedness of direct regression by modeling the output distribution rather than predicting point estimates, narrowing the performance gap with heatmap methods.

A recent breakthrough that elegantly bridges the gap between heatmaps and regression is SimDR (Simple Disentangled Representation), presented by Li~\emph{et al.}~\cite{li20212d} as an oral paper at ECCV 2021. SimDR poses the provocative question: \emph{is a 2D heatmap representation even necessary for accurate keypoint localization?} Their answer reformulates pose estimation as two independent 1D classification problems---predicting the horizontal bin index and the vertical bin index for each keypoint---thereby replacing the quadratic-cost 2D heatmap ($H \times W$ per keypoint) with two linear-cost 1D probability vectors ($H + W$ per keypoint). By using a configurable bin resolution (scale factor $s$) that exceeds the spatial resolution of the feature map, SimDR achieves sub-pixel accuracy while entirely avoiding 2D spatial quantization error. The 1D classification targets are soft Gaussian distributions supervised by KL divergence loss, providing a smoother optimization landscape than the MSE loss typically used for 2D heatmaps.

In the proposed framework, the SimDR formulation is applied to $48 \times 48$ RoI features with scale factor $s = 2$ (yielding 96 bins per axis), enabling sub-pixel keypoint localization within the extremely limited spatial resolution of the thermal input images. An auxiliary visibility branch predicts per-keypoint occlusion flags via binary cross-entropy loss, following the three-level visibility weighting scheme of CrowdPose~\cite{li2019crowdpose} (invisible $\rightarrow$ ignored; occluded $\rightarrow$ reduced weight; visible $\rightarrow$ full weight). Through systematic ablation experiments (Chapter Experiments), this approach is compared against four alternative head designs---standard 2D Gaussian heatmap regression, foreground-weighted heatmap with visibility prediction, soft-argmax integral regression, and direct coordinate regression---demonstrating that SimDR achieves the best accuracy--efficiency trade-off for the low-resolution thermal input scenario.

\section{Skeleton-Based Action Recognition}
\label{sec:rw_action}

\subsection*{Motivation and Advantages}

Skeleton-based action recognition has emerged as a prominent research direction that leverages the compact, structured representation of human body keypoints for temporal activity understanding. Compared to appearance-based methods such as two-stream networks and 3D convolutional networks that operate on raw pixel intensities or optical flow, skeleton representations offer several compelling advantages. They are inherently robust to background clutter, illumination changes, and appearance variation, since only the geometric configuration of body joints is retained. They are also highly compact---a skeleton with $K$ keypoints requires only $2K$ or $3K$ values per frame (2D or 3D coordinates), compared to millions of pixel values in a video frame---making skeleton-based methods computationally efficient and well suited for edge deployment. Furthermore, in privacy-sensitive thermal infrared applications where texture and color information are absent, the skeleton representation naturally aligns with the available information: the thermal image primarily reveals the spatial extent and configuration of the body, which is precisely what keypoints encode.

\subsection*{RNN and LSTM Approaches}

Early deep learning approaches to skeleton-based action recognition employed recurrent neural networks (RNNs) to model the temporal evolution of keypoint sequences. Du~\emph{et al.}~\cite{du2015hierarchical} proposed a hierarchical RNN architecture that decomposes the skeleton into five anatomical sub-groups (two arms, two legs, and trunk), processes each sub-group with a separate RNN branch, and hierarchically fuses the temporal features in a tree structure. This decomposition exploits the natural hierarchy of the human body and demonstrated the viability of sequential modeling for activity classification on early skeleton benchmarks. The Gated Recurrent Unit (GRU)~\cite{cho2014learning} and Long Short-Term Memory (LSTM) networks were subsequently applied across various skeleton benchmarks, offering mechanisms to capture long-range temporal dependencies through learned gating functions that control information flow. These recurrent approaches established strong baselines, particularly for datasets with relatively few keypoints and simple skeleton topologies.

\subsection*{Graph Convolutional Networks}

The introduction of Spatial-Temporal Graph Convolutional Networks (ST-GCN) by Yan~\emph{et al.}~\cite{yan2018spatial} marked a paradigm shift in the field. ST-GCN formulates the skeleton as a spatiotemporal graph in which joints are nodes and bones (anatomical connections) are edges, and applies graph convolutions that simultaneously capture spatial dependencies among joints within a single frame and temporal dynamics across consecutive frames. By operating directly on the graph structure, ST-GCN avoids the need to flatten the skeleton into a 1D vector (as required by RNNs) and can naturally model the relational structure between body parts. ST-GCN achieved state-of-the-art results on the NTU~RGB+D dataset~\cite{shahroudy2016ntu}, one of the largest benchmarks for skeleton-based action recognition, and on the Kinetics skeleton dataset.

Subsequent works extended the ST-GCN paradigm in several important directions. Two-Stream Adaptive Graph Convolutional Networks (2s-AGCN)~\cite{shi2019two} introduced two key innovations: (1) learnable adjacency matrices that adapt the graph topology during training, allowing the network to discover task-relevant joint relationships beyond the predefined anatomical connections; and (2) a two-stream architecture that processes both joint position data and bone vector data (computed as difference vectors between connected joints) in parallel, fusing the two streams for final classification. The bone stream provides complementary information about limb orientations and lengths that is not directly encoded in isolated joint positions. MS-G3D~\cite{liu2020msg3d} further improved multi-scale temporal modeling through disentangled multi-scale graph convolutions that capture temporal dependencies at different time scales simultaneously, addressing the limitation of ST-GCN's fixed temporal receptive field. Channel-wise Topology Refinement Graph Convolution (CTR-GC)~\cite{chen2021channel} refined the graph topology at the channel level, learning different relational structures for different feature channels and achieving strong performance on NTU~RGB+D benchmarks.

\subsection*{Input Representations and Multi-Stream Fusion}

Beyond architectural innovations, the choice of input representation has proven critical for skeleton-based recognition performance. Raw joint coordinates provide the basic positional information but are sensitive to the person's absolute position and scale within the frame. Bone features, computed as the difference vectors between connected joint pairs, encode limb orientation and length information that is invariant to translation. Joint velocity features, obtained by computing first-order temporal differences of joint positions between consecutive frames, capture motion dynamics and are particularly informative for distinguishing actions that share similar static poses but differ in movement patterns (e.g., standing still vs.\ standing up). Bone velocity features combine both geometric and temporal information. Multi-stream fusion of these four input modalities---joint positions, bone vectors, joint velocities, and bone velocities---has become standard practice in high-performing systems~\cite{shi2019two, chen2021channel}, with each stream processed by a separate network and the outputs fused (typically by score averaging) for final prediction.

The NTU~RGB+D dataset~\cite{shahroudy2016ntu}, captured using Microsoft Kinect depth sensors with 25 body keypoints per person, has served as the primary benchmark for evaluating skeleton-based methods. However, the 25-keypoint Kinect skeleton is substantially richer than the 7-keypoint skeleton available in the thermal infrared setup of this work, and reduced skeleton complexity presents unique challenges for recognition algorithms that rely on rich inter-joint spatial structure.

\subsection*{The Proposed Approach}

The proposed framework adopts a lightweight GRU-based temporal classifier rather than a full graph convolutional network. This design choice is motivated by the observation that the skeleton in this work consists of only seven keypoints (head, left/right shoulders, left/right hands, hip, and lower body reference) with a simple, largely linear topology, for which the expressive power of graph convolutions offers diminishing returns relative to their computational overhead. For a 7-joint skeleton, the spatial graph has very limited structure---far less than the 25-joint Kinect skeleton for which ST-GCN was designed---and a GRU operating on the flattened joint feature vector can effectively capture the relevant temporal patterns. The feature representation evolves through several design iterations during development: from raw normalized keypoint coordinates, to velocity-augmented features that capture inter-frame motion dynamics, to bone-angle features that encode geometric relationships between limbs, and finally to a variant that fuses skeleton features with appearance-level RoI embeddings extracted from the detection backbone. This multi-task learning approach~\cite{caruana1997multitask} allows the action classifier to benefit from shared feature representations learned jointly with the detection and pose estimation tasks.

\section{Fall Detection Systems}
\label{sec:rw_fall}

\subsection*{Sensing Modalities}

Fall detection has been the subject of extensive research over the past two decades, driven by the significant health consequences of undetected falls among the elderly~\cite{who2007falls, tinetti1988risk}. Falls are a leading cause of injury-related morbidity and mortality among older adults, and the duration of time spent on the ground after a fall---the so-called ``long lie''---is strongly correlated with adverse outcomes including hypothermia, dehydration, and death~\cite{lord2007falls}. With the global population aging rapidly, the demand for reliable, automated fall detection systems in indoor environments has become increasingly urgent. Comprehensive surveys of fall detection technologies have been provided by Mubashir~\emph{et al.}~\cite{mubashir2013survey}, Wang~\emph{et al.}~\cite{wang2020elderly}, and Igual~\emph{et al.}~\cite{igual2013challenges}, each offering different taxonomies and evaluation perspectives.

Existing fall detection systems can be categorized by their sensing modality into three broad classes: wearable sensors, ambient sensors, and vision-based systems. Wearable sensor-based approaches employ tri-axial accelerometers, gyroscopes, or inertial measurement units (IMUs) attached to the body (typically at the waist, wrist, or chest) to detect falls by analyzing abrupt changes in acceleration magnitude, angular velocity, and body orientation~\cite{mubashir2013survey}. While these methods can achieve high sensitivity and specificity under controlled conditions, they fundamentally depend on user compliance: the individual must consistently wear, charge, and maintain the device. Studies have reported low long-term adherence rates among elderly populations, particularly those with cognitive impairments, severely limiting the practical utility of wearable solutions in unsupervised real-world settings~\cite{lord2007falls}.

Ambient sensor-based approaches offer a non-intrusive alternative by instrumenting the environment rather than the person. Floor-mounted pressure or vibration sensors can detect the impact of a fall event, while acoustic sensors can capture characteristic sound patterns associated with falls. However, these methods are limited in spatial coverage, lack contextual information about the nature of the event, and are prone to false alarms in realistic home environments.

\subsection*{Vision-Based Fall Detection}

Vision-based fall detection has attracted growing interest due to its ability to monitor large areas continuously without requiring user cooperation or environmental instrumentation. Early approaches using RGB cameras relied on handcrafted features extracted from video sequences. Rougier~\emph{et al.}~\cite{rougier2011robust} proposed a fall detection method based on tracking changes in the human silhouette shape over time, using Gaussian mixture models for foreground segmentation and shape matching for event classification. Charfi~\emph{et al.}~\cite{charfi2013optimized} systematically evaluated spatio-temporal descriptors for real-time fall detection, comparing SVM and AdaBoost classifiers on features including motion history images, aspect ratio trajectories, and velocity profiles. These handcrafted approaches established important baselines but were limited by their reliance on manually designed features that may not generalize across diverse environments and fall patterns.

The application of deep learning to vision-based fall detection brought substantial performance improvements, with CNNs learning discriminative spatio-temporal representations directly from video data. Depth cameras (e.g., Microsoft Kinect) provided an intermediate solution that captures 3D body structure while partially preserving privacy through the absence of color and texture information~\cite{mastorakis2014fall}. However, depth sensors have limited operational range (typically 0.5--4.5~m for Kinect), are susceptible to interference from reflective and transparent surfaces, and still capture sufficient geometric detail to potentially identify individuals through body shape analysis.

\subsection*{Thermal Infrared Fall Detection}

Thermal infrared cameras have emerged as the most compelling modality for fall detection in privacy-critical environments, combining the spatial awareness of vision-based systems with the anonymity guarantees required by data protection regulations. Several research groups have demonstrated the viability of thermal fall detection at various sensor resolutions.

At the higher end of the resolution spectrum, Vadivelu~\emph{et al.}~\cite{vadivelu2016thermal} explored thermal imaging for elderly fall detection using standard thermal cameras, demonstrating robust performance across varying lighting conditions. Tao~\emph{et al.}~\cite{tao2018real} proposed a deep learning pipeline for real-time human body segmentation and fall detection from thermal infrared images, combining semantic segmentation with temporal analysis of body orientation changes. Nogas~\emph{et al.}~\cite{nogas2020fall} demonstrated that CNNs trained on thermal images from ceiling-mounted cameras can reliably detect falls even at relatively low resolutions, achieving strong performance with compact network architectures suitable for embedded deployment.

At the ultra-low-resolution end---which offers the strongest privacy guarantees---Trofimova~\emph{et al.}~\cite{trofimova2017indoor} investigated activity recognition using low-resolution infrared array sensors, showing that basic activities including sitting, standing, lying, and walking can be distinguished even from coarse thermal signatures. Tateno~\emph{et al.}~\cite{tateno2020privacy} specifically addressed privacy-preserving fall detection in smart homes using low-resolution thermal array sensors, proposing a CNN-based classification approach that achieves reliable fall detection while ensuring that the spatial resolution is too low to reveal any personally identifiable information. The present work builds upon this body of research by operating on $40 \times 114$ pixel thermal images produced by dual Heimann HTPA sensors~\cite{heimann2024htpa}, a resolution that provides sufficient spatial detail for 7-keypoint pose estimation while remaining firmly within the privacy-preserving regime.

\subsection*{Binary versus Multi-Class Formulation}

A fundamental design question in fall detection is the choice of problem formulation: binary classification (fall vs.\ non-fall) versus multi-class action recognition from which falls are inferred indirectly~\cite{wang2020elderly}. The binary formulation directly optimizes for the target event by training a classifier specifically to distinguish falls from all other activities. This approach avoids error propagation from intermediate action labels and can concentrate model capacity on the decision boundary that matters most, but it may miss contextual information about the transition dynamics leading up to a fall (e.g., the distinction between falling while walking versus falling while trying to stand). Multi-class approaches~\cite{zhang2012rgb} can capture finer activity semantics and provide richer behavioral analytics, but they require more extensive annotation, increase the number of output classes over which the classifier must discriminate, and risk diluting the classifier's sensitivity to the relatively rare fall class---a concern exacerbated by the severe class imbalance inherent in realistic fall datasets, where falls constitute only a small fraction of all observed activities.

Rule-based post-processing heuristics have been employed in both paradigms to reduce false positive rates. Common heuristics include detecting rapid downward displacement of the head keypoint (a kinematic signature of falling), sustained horizontal body orientation following a motion event (distinguishing falls from intentionally lying down), prolonged stillness in a lying-down posture (detecting ``long lie'' events), and sudden changes in bounding box aspect ratio from vertical to horizontal~\cite{charfi2013optimized, rougier2011robust}. These rule-based methods complement learned classifiers by encoding domain knowledge about the biomechanics and temporal dynamics of fall events.

\subsection*{The Proposed Approach}

In the proposed system, fall detection is formulated as a binary temporal classification task, directly predicting a fall probability from a sequence of normalized skeleton keypoints using a GRU-based classifier. Through ablation experiments in Chapter Experiments, it is demonstrated that this direct binary formulation outperforms indirect inference from the 10-class activity taxonomy available in the SenIR dataset when evaluated on fall-specific metrics (precision, recall, and F1-score). To address the inherent class imbalance---fall events are substantially rarer than normal activities in any realistic monitoring scenario---a weighted binary cross-entropy loss with positive class weighting ($w_+ = 5.0$) is employed, following standard practices for imbalanced classification in medical and safety-critical domains.

The temporal integration facilitated by ByteTrack~\cite{zhang2022bytetrack} multi-object tracking, built upon the Kalman Filter motion model of SORT~\cite{bewley2016simple}, enables the classifier to operate over variable-length keypoint sequences (1--30 frames). This temporal context allows the GRU to observe the dynamic transition pattern---the rapid downward displacement of the head, the rotation of the body axis from vertical to horizontal, the deceleration following impact---that distinguishes genuine falls from similar static configurations such as a person lying in bed. The end-to-end integration of detection, pose estimation, action classification, and tracking within a single jointly trained framework is a distinguishing feature of the proposed approach, enabling shared feature representations and efficient inference on the resource-constrained edge hardware targeted by the SensIR project.